\documentclass[11pt]{article}
\usepackage[a4paper,margin=27mm]{geometry}
\usepackage[T1]{fontenc}
\usepackage{lmodern,microtype,amsmath,amssymb,amsthm,mathtools}
\usepackage{booktabs,longtable,array,enumitem,xcolor,xurl}
\usepackage[colorlinks=true,linkcolor=blue!45!black,citecolor=blue!45!black,urlcolor=blue!45!black]{hyperref}
\usepackage{bookmark}
\setlength{\parskip}{4pt}
\setlength{\parindent}{15pt}
\setlength{\emergencystretch}{3em}
\numberwithin{equation}{section}
\newtheorem{theorem}{Theorem}[section]
\newtheorem{proposition}[theorem]{Proposition}
\newtheorem{lemma}[theorem]{Lemma}
\newtheorem{corollary}[theorem]{Corollary}
\theoremstyle{definition}
\newtheorem{definition}[theorem]{Definition}
\theoremstyle{remark}
\newtheorem{remark}[theorem]{Remark}
\newtheorem{example}[theorem]{Example}
\DeclareMathOperator{\Li}{Li}
\DeclareMathOperator{\Fp}{Fp}
\DeclareMathOperator{\Rea}{Re}
\DeclareMathOperator{\CT}{CT}
\newcommand{\T}{\mathbb T}
\newcommand{\C}{\mathbb C}
\newcommand{\Q}{\mathbb Q}
\newcommand{\Z}{\mathbb Z}
\newcommand{\N}{\mathbb N}
\newcommand{\Id}{\mathsf I}
\newcommand{\G}{\mathcal G}
\newcommand{\dd}{\,\mathrm d}
\newcommand{\cB}{\mathcal B}
\newcommand{\zj}[3]{\zeta^{(#1)}(#2,#3)}
\newcommand{\pin}{\texttt{0efdbaa6944357e578a8839ec88d040c6e05681c}}
\newcommand{\Liord}[2]{\Li_{#1}^{[#2]}}
\title{A Convolution Algebra of Stieltjes Functions\\[4pt]
\large Finite-part identities, polygamma contact terms,\\log-gamma convolution powers, and polylogarithmic correlations}
\author{Research continuation prepared for Vladimir Reshetnikov\\\small ProveIt / Analysis / Polylogarithms}
\date{10 October 2026}
\hypersetup{pdftitle={A Convolution Algebra of Stieltjes Functions},pdfauthor={Research continuation prepared for Vladimir Reshetnikov},pdfsubject={Exact identities and normalization-aware Hurwitz-jet calculus}}
\begin{document}
\maketitle
\begin{abstract}
The Hurwitz-zeta convolution semigroup has a singular parameter at which
its scalar formula ceases to describe the full periodic distribution. Keeping
the missing contact term produces a closed convolution algebra for all
generalized Stieltjes functions. We prove an all-index Bell-polynomial
normal form, a bivariate Gamma-quotient multiplication law, and absolutely
convergent subtracted-integral realizations away from the singular point.
The first identity is a finite-part digamma convolution equal to
$2\gamma_1(x)+\zeta(2)$. Differentiation gives an explicit anomaly:
periodic differentiation and canonical Hadamard finite part differ by a
known multiple of a derivative of the delta distribution. Consequently,
every convolution of two finite-part polygamma functions reduces to just
two Hurwitz-zeta jets, with corrected harmonic-number coefficients.
At the integrated end, every circular convolution power of centered
$\log\Gamma$ is a finite combination of jets at a single negative integer.
Reflection changes the same spectral calculation into order derivatives
of polylogarithms and yields a full shifted log-gamma correlation formula.
All analytic statements are proved in the article. Exact finite algebra
and independent high-precision numerical diagnostics are supplied separately.
The basic nonsingular semigroup is attributed to prior work; no priority
claim, period-independence claim, or proof-assistant verification is made.
\end{abstract}

\clearpage
\tableofcontents
\clearpage

\section{Scope, provenance, and main conclusions}
\label{sec:scope}
This continuation addresses identities rather than inequalities. Its
starting point is the integrated and differentiated Hurwitz-jet calculus
in the ProveIt manuscript~\cite{ProveIt}, especially the chapters
\texttt{07-integration.tex} and \texttt{08-differentiation.tex}.
The repository snapshot used for the final source audit is
\begin{center}\small\pin.\end{center}
The surviving-problems chapter in that snapshot already calls the
$S_4$ identity proved. We do not present it as an open problem, nor claim a
new solution of the remaining $S_6$ or arithmetic-independence questions.
The incoming-directory inventory and intake instructions were inspected;
the binary incoming archives were not individually extracted or
exhaustively audited. These limits matter when comparing this contribution
with a rapidly changing research repository. Detailed source paths and
suggested placement are in the accompanying audit and integration notes.

\subsection{What is established here}
The principal addition is a normalization-aware identity calculus. If
$G_n$ is the canonical periodic finite part of $\gamma_n$ and
$\Id=\delta_0-1$, then
\begin{equation}\label{eq:intro-first}
 G_0*G_0=2G_1-\zeta(2)\Id.
\end{equation}
The star denotes circular convolution, not pointwise multiplication.
Restricting away from $0\pmod1$ changes the last term into $+\zeta(2)$.
More generally, all $G_m*G_n$ are finite linear combinations of
$\Id,G_0,\ldots,G_{m+n+1}$, with explicit zeta-polynomial coefficients.
There is no unresolved analytic remainder.

The derivative correction is equally explicit. Write
$e_j(k)=e_j(1,1/2,\ldots,1/k)$, setting $e_j(k)=0$ for $j>k$. Then
\begin{equation}\label{eq:intro-contact}
 D^kG_n=\Fp_{\T}[\gamma_n^{(k)}]
       +(-1)^{n+1}n!e_{n+1}(k)\delta_0^{(k)}\qquad(k\ge1).
\end{equation}
The ordinary pointwise parameter-derivative theorem in the repository
remains valid. Equation~\eqref{eq:intro-contact} specifies the extra term
needed when that theorem is extended to periodic finite parts.

For example, putting $\Psi_k=\Fp_{\T}\psi^{(k)}$, one obtains, for $0<x<1$,
\begin{equation}\label{eq:intro-tri}
 (\Psi_1*\Psi_1)(x)=-4\zeta'(3,x)-2\zeta(3,x).
\end{equation}
A completely explicit, absolutely convergent subtracted-integral version
appears in Section~\ref{sec:poly}. Thus distribution notation is a tool
for derivation, not an obstacle to ordinary numerical evaluation.

\subsection{Classical input and novelty boundaries}
Hurwitz's Fourier formula and the usual Gamma, polygamma, and polylogarithm
identities are classical; convenient primary reference compilations are
DLMF Sections 25.11 and 25.12~\cite{DLMFHz,DLMFLi}. The nonsingular
Hurwitz convolution law below is already stated by Zhi-Hong Sun in
Remark 3.2 of \emph{On the properties of invariant functions}~\cite{Sun},
for real parameters greater than one. We supply a short Fourier proof and
continue it as an entire distribution-valued family. Espinosa and Moll's
integral calculus~\cite{EM1,EM2} supplies important context at the integrated
end; Coffey's all-order Stieltjes derivative formulas~\cite{Coffey} supply
context at the differentiated end.

The all-index finite-part multiplication table, explicit local
subtractions, contact correction, polygamma collapse, and convolution-power
formulas are the concrete results developed in this continuation. A
literature search located the basic semigroup and the classical moment
and derivative theories, but did not establish a comprehensive priority
history for every specialization below. Accordingly, ``proved here'' is
not used as a synonym for ``first proved in the literature.''

\section{Fourier conventions and the Hurwitz semigroup}
\label{sec:semigroup}
We work on $\T=\mathbb R/\mathbb Z$, with normalized Lebesgue measure.
For a periodic distribution $T$, use the convention
\[
 \widehat T(k)=\langle T,e^{-2\pi i kx}\rangle,\qquad k\in\Z.
\]
The constant function $1$ and the point mass $\delta_0$ both have mean
one, but their nonzero Fourier coefficients are respectively $0$ and $1$.
Thus
\[
 \Id=\delta_0-1
\]
is the convolution identity on the space of mean-zero distributions. In
particular, $\Id*T=T$ whenever $\widehat T(0)=0$.

The logarithm used throughout is
\begin{equation}\label{eq:logbranch}
 \lambda_k=\log(2\pi i k)
 =\log(2\pi|k|)+\frac{\pi i}{2}\operatorname{sgn}k,
 \qquad k\ne0.
\end{equation}
All powers below refer to this fixed branch.

\begin{definition}\label{def:R}
For $\alpha\in\C$, define $R_\alpha$ by
\begin{equation}\label{eq:Rfourier}
 \widehat R_\alpha(0)=0,\qquad
 \widehat R_\alpha(k)=e^{-\alpha\lambda_k}\quad(k\ne0).
\end{equation}
Equivalently, its Fourier series is
$\sum_{k\ne0}(2\pi i k)^{-\alpha}e^{2\pi i kx}$ in the distribution sense.
\end{definition}

\begin{lemma}[Entire distribution family]\label{lem:entire}
The family $\alpha\mapsto R_\alpha$ is entire as a periodic-distribution
family. It is represented by an absolutely convergent Fourier series for
$\Rea\alpha>1$. For $\Rea\alpha>0$, it is the locally integrable periodic
extension of
\begin{equation}\label{eq:Rhurwitz}
 R_\alpha(x)=\frac{\zeta(1-\alpha,x)}{\Gamma(\alpha)},\qquad0<x<1.
\end{equation}
\end{lemma}
\begin{proof}
On each compact set of parameters, the coefficients in
\eqref{eq:Rfourier}, together with any fixed number of parameter derivatives,
have at most polynomial growth in $|k|$. Fourier coefficients of a smooth
test function decay faster than every polynomial. Hence the paired Fourier
series and all its parameter derivatives converge locally uniformly. This
proves the first assertion; the absolute-convergence assertion follows
immediately from $\sum k^{-\Rea\alpha}$.

For $\Rea\alpha>1$, equation~\eqref{eq:Rhurwitz} is precisely Hurwitz's
Fourier formula. To continue to $\Rea\alpha>0$, use
\[
 \zeta(1-\alpha,x)=x^{\alpha-1}+\zeta(1-\alpha,1+x).
\]
On compact parameter subsets of that half-plane, the first term and its
parameter derivatives are dominated by integrable powers of $x$ times
powers of $|\log x|$. The second term is regular in $x\in[0,1]$; any
apparent parameter singularity at $\alpha=0$ is outside the stated open
half-plane. The two holomorphic distribution families agree in an open
subregion, and hence agree throughout the half-plane.
\end{proof}

\begin{theorem}[Semigroup and differential ladder]\label{thm:semigroup}
For all $\alpha,\beta\in\C$,
\begin{equation}\label{eq:semigroup}
 R_\alpha*R_\beta=R_{\alpha+\beta},\qquad
 DR_\alpha=R_{\alpha-1},\qquad R_0=\Id.
\end{equation}
If $J_{\alpha,r}=\partial_\alpha^rR_\alpha$, then
\begin{equation}\label{eq:mixedjets}
 J_{\alpha,r}*J_{\beta,s}=J_{\alpha+\beta,r+s},\qquad
 D^kJ_{\alpha,r}=J_{\alpha-k,r}.
\end{equation}
\end{theorem}
\begin{proof}
Periodic distributions can be convolved: equivalently, their
polynomially growing Fourier coefficients can be multiplied. For $k\ne0$,
\[
 e^{-\alpha\lambda_k}e^{-\beta\lambda_k}
 =e^{-(\alpha+\beta)\lambda_k},\qquad
 (2\pi i k)e^{-\alpha\lambda_k}=e^{-(\alpha-1)\lambda_k}.
\]
Both sides have zero zeroth coefficient. Fourier uniqueness proves
\eqref{eq:semigroup}. Differentiating the locally uniformly paired
Fourier series gives~\eqref{eq:mixedjets}.
\end{proof}

For real $\alpha,\beta>1$, the first identity, expressed using
\eqref{eq:Rhurwitz}, is the nonsingular law in Sun~\cite[Remark 3.2]{Sun}.
The new use here is at its singular integer parameters and their spectral
jets. At positive integers,
\[
 R_m(x)=-\frac{B_m(\{x\})}{m!}\quad(m\ge1),
\]
up to immaterial point values of the sawtooth for $m=1$. At negative
integers,
\begin{equation}\label{eq:negativeR}
 R_{-k}=\delta_0^{(k)}\quad(k\ge1).
\end{equation}
The quotient in~\eqref{eq:Rhurwitz} vanishes pointwise at $\alpha=-k$
away from the integers, while the distribution is nonzero. This is the
basic reason contact terms must be retained.

\section{Canonical finite parts of the Stieltjes functions}
\label{sec:finiteparts}
The generalized Stieltjes functions are defined for $x>0$ by
\begin{equation}\label{eq:laurent}
 \zeta(1-z,x)=-\frac1z+\sum_{n=0}^\infty\gamma_n(x)\frac{z^n}{n!}.
\end{equation}
Thus $\gamma_0(x)=-\psi(x)$. The Hurwitz shift relation gives
\begin{equation}\label{eq:stshift}
 \gamma_n(x)=\frac{\log^n x}{x}+\gamma_n(1+x).
\end{equation}
Consequently, periodic extension of $\gamma_n$ is not locally integrable
at $0$; a convention is necessary.

\begin{definition}[Canonical periodic finite part]\label{def:G}
For a smooth periodic test function $\phi$, define
\begin{align}
 \langle G_n,\phi\rangle
 &=\int_0^1\left(\gamma_n(x)\phi(x)
       -\frac{\log^n x}{x}\phi(0)\right)\dd x\label{eq:Gdefinition}\\
 &=\lim_{\epsilon\downarrow0}\left\{
       \int_\epsilon^1\gamma_n(x)\phi(x)\dd x
       +\frac{\log^{n+1}\epsilon}{n+1}\phi(0)\right\}.
       \label{eq:Gcutoff}
\end{align}
\end{definition}
The integrand in~\eqref{eq:Gdefinition} is integrable: its singular part is
$\log^n x\,[\phi(x)-\phi(0)]/x$, which is $O(|\log x|^n)$.
No arbitrary finite multiple of $\delta_0$ is added.

\begin{lemma}[The distributional Laurent expansion]\label{lem:Laurentdist}
Let
\[
 Z(s)=\Gamma(1-s)R_{1-s}.
\]
It is a meromorphic periodic-distribution family, agreeing with the
periodic locally integrable Hurwitz zeta for $\Rea s<1$. Near $z=0$,
\begin{equation}\label{eq:ZLaurent}
 Z(1-z)=\frac{\Id}{z}+\sum_{n=0}^\infty G_n\frac{z^n}{n!}.
\end{equation}
Every $G_n$ has mean zero.
\end{lemma}
\begin{proof}
The meromorphic continuation follows from Lemma~\ref{lem:entire}. For
$\Rea z>0$, split the pairing using
$\zeta(1-z,x)=x^{z-1}+\zeta(1-z,1+x)$. The first term has expansion
\[
 \int_0^1x^{z-1}\phi(x)\dd x
 =\frac{\phi(0)}z+
   \sum_{n\ge0}\frac{z^n}{n!}
   \int_0^1\frac{\log^n x}{x}[\phi(x)-\phi(0)]\dd x.
\]
The second term contributes $-z^{-1}\int_0^1\phi$ and the coefficients
$\int_0^1\gamma_n(1+x)\phi(x)\dd x$. These expansions converge locally
near zero after pairing, by the same endpoint domination used above.
Combining them proves~\eqref{eq:ZLaurent} with the exact subtraction in
Definition~\ref{def:G}. Since $\widehat Z(s)(0)=0$, every coefficient has
mean zero. Equivalently, direct integration over $[1,2]$ gives
$\int_1^2\zeta(1-z,x)\dd x=-1/z$, whose regular coefficients vanish.
\end{proof}

\begin{remark}[Two different residues]
The scalar function in~\eqref{eq:laurent} has residue $-1$ as a function
of $z$, whereas the periodic-distribution family has residue
$\delta_0-1$. Analytic continuation across a nonintegrable endpoint adds
a point-supported residue. Treating the scalar residue as the whole
periodic residue gives wrong convolution constants.
\end{remark}

\section{All-index convolution closure and exact normal forms}
\label{sec:algebra}
Write $\G(z)=\sum_{n\ge0}G_nz^n/n!$, holomorphic as a distribution family
for $|z|<1$. By~\eqref{eq:ZLaurent},
\begin{equation}\label{eq:masterG}
 \Id+z\G(z)=\Gamma(1+z)R_z.
\end{equation}
Let $X=G_0$. For a complete exponential Bell polynomial we use
\begin{equation}\label{eq:Belldef}
 \exp\left(\sum_{j\ge1}b_j\frac{z^j}{j!}\right)
 =\sum_{r\ge0}\cB_r(b_1,\ldots,b_r)\frac{z^r}{r!}.
\end{equation}
When its arguments are distributions, all products in this definition
are convolution products and its scalar unit is $\Id$.

\begin{theorem}[One-generator convolution algebra]\label{thm:Bell}
For every $n\ge0$,
\begin{equation}\label{eq:BellG}
 (n+1)G_n=\cB_{n+1}(b_1,\ldots,b_{n+1}),\qquad
 b_1=X,\quad b_j=(-1)^j(j-1)!\zeta(j)\Id\ (j\ge2).
\end{equation}
The span of $\Id,G_0,G_1,\ldots$ is a convolution algebra isomorphic to
the polynomial algebra $\C[X]$. These distributions form a basis of that
algebra, with the unit listed first.
\end{theorem}
\begin{proof}
For $k\ne0$, the Fourier transform of~\eqref{eq:masterG} is
\[
 1+z\widehat\G(z)(k)
 =\Gamma(1+z)e^{-z\lambda_k}
 =\exp\left[(-\gamma-\lambda_k)z+
       \sum_{j\ge2}\frac{(-1)^j\zeta(j)}jz^j\right].
\]
The linear coefficient is
$\widehat X(k)=-\gamma-\lambda_k$. Comparing with
\eqref{eq:Belldef} proves~\eqref{eq:BellG}; the zero Fourier coefficient
of each distributional expression vanishes. The leading term in $G_n$
is $X^{*(n+1)}/(n+1)$, so the listed distributions span exactly the
polynomial algebra generated by $X$.

There are no polynomial relations among its powers. Indeed, if
$P(X)=0$ distributionally, then
$P(-\gamma-\log(2\pi k)-i\pi/2)=0$ for every positive integer $k$.
These are infinitely many distinct complex numbers, and hence $P$ is the
zero polynomial. This also proves the basis assertion.
\end{proof}

This algebraic freeness concerns distributions as functions of the
variable $x$. It is not an independence theorem about their values at a
fixed rational point, nor about zeta or Stieltjes constants as numbers.

The first Bell forms are
\begin{align}
 2G_1&=X^{*2}+\zeta(2)\Id,\label{eq:Belllow}\\
 3G_2&=X^{*3}+3\zeta(2)X-2\zeta(3)\Id,\notag\\
 4G_3&=X^{*4}+6\zeta(2)X^{*2}-8\zeta(3)X
                +[3\zeta(2)^2+6\zeta(4)]\Id.\notag
\end{align}
They already imply many exact integrals. An efficient all-index
multiplication rule is supplied by the following Gamma quotient.

\begin{theorem}[Bivariate multiplication law]\label{thm:product}
Define, near $(u,v)=(0,0)$,
\begin{align}
 B(u,v)&=\frac{\Gamma(1+u)\Gamma(1+v)}{\Gamma(1+u+v)}\label{eq:Buv}\\
 &=\exp\left(\sum_{j\ge2}\frac{(-1)^j\zeta(j)}j
          [u^j+v^j-(u+v)^j]\right).\notag
\end{align}
Then
\begin{align}
 uv\,\G(u)*\G(v)={}&[B(u,v)-1]\Id\label{eq:Gproduct}\\
 &+B(u,v)(u+v)\G(u+v)-u\G(u)-v\G(v).\notag
\end{align}
The right side is divisible by $uv$ as a convergent distribution-valued
power series. In particular,
\begin{equation}\label{eq:topindex}
 G_m*G_n=\frac{m+n+2}{(m+1)(n+1)}G_{m+n+1}
       +\operatorname{span}_{\C}\{\Id,G_0,\ldots,G_{m+n}\},
\end{equation}
and every coefficient is a polynomial over $\Q$ in zeta values at positive
integers. Its complete value is obtained by multiplying the coefficient
of $u^{m+1}v^{n+1}$ in~\eqref{eq:Gproduct} by $m!n!$.
\end{theorem}
\begin{proof}
Convolve~\eqref{eq:masterG} at parameters $u$ and $v$ and apply
Theorem~\ref{thm:semigroup}. The result is
\[
 [\Id+u\G(u)]*[\Id+v\G(v)]
 =B(u,v)[\Id+(u+v)\G(u+v)].
\]
Subtract the unit and the two linear summands. Both restrictions $u=0$
and $v=0$ vanish identically, proving divisibility. Coefficient extraction
is legitimate in a sufficiently small bidisc. The largest-index term
comes from $(u+v)\G(u+v)$ with the constant coefficient of $B$; its
coefficient is
$\frac{m!n!}{(m+n+1)!}\binom{m+n+2}{m+1}$, which equals the number in
\eqref{eq:topindex}. The exponential expansion of $B$ proves the
zeta-polynomial assertion.
\end{proof}

Four useful multiplication rules are
\begin{align}
 G_0*G_0&=2G_1-\zeta(2)\Id,\label{eq:product00}\\
 G_0*G_1&=\frac32G_2-\zeta(2)G_0+\zeta(3)\Id,\label{eq:product01}\\
 G_0*G_2&=\frac43G_3-2\zeta(2)G_1+2\zeta(3)G_0-2\zeta(4)\Id,\label{eq:product02}\\
 G_1*G_1&=G_3-2\zeta(2)G_1+2\zeta(3)G_0-\frac14\zeta(4)\Id.
 \label{eq:product11}
\end{align}
The last simplification uses $\zeta(2)^2=5\zeta(4)/2$. The verification
catalog retains independent formal symbols for zeta values until such a
specific relation is explicitly applied.

\section{Ordinary integral realizations and explicit Stieltjes identities}
\label{sec:integrals}
The singular support of $G_n$ is contained in $\{0\}$. Thus its
convolution with another $G_m$ is smooth away from $0$: in a local
convolution integral at $0<x<1$, the two singular factors occur at the
distinct integration points $t=0$ and $t=x$. We now give a direct formula
that also proves this local realization without requiring a product of
distributions at the same point.

\begin{theorem}[Absolutely convergent subtracted integrals]\label{thm:integral}
Put $a_n(t)=\log^n(t)/t$. For $m,n\ge0$ and $0<x<1$, define
\begin{align}
 C_{m,n}(x)={}&\int_0^x\big[\gamma_m(t)\gamma_n(x-t)
          -\gamma_n(x)a_m(t)-\gamma_m(x)a_n(x-t)\big]\dd t
          \label{eq:Cmn}\\
 &+\int_x^1\gamma_m(t)\gamma_n(1+x-t)\dd t\notag\\
 &+\gamma_n(x)\frac{\log^{m+1}x}{m+1}
  +\gamma_m(x)\frac{\log^{n+1}x}{n+1}.\notag
\end{align}
Both displayed integrals are absolutely convergent, and
$C_{m,n}(x)=(G_m*G_n)(x)$. In particular, every expression
\eqref{eq:Cmn} has a finite exact reduction using Theorem~\ref{thm:product}.
\end{theorem}
\begin{proof}
Near $t=0$, the shift relation~\eqref{eq:stshift} leaves a possible
singularity
\[
 a_m(t)[\gamma_n(x-t)-\gamma_n(x)]=O(|\log t|^m).
\]
The other endpoint is the same argument with $m,n$ interchanged. The
second integral is nonsingular. These observations prove absolute
convergence, locally uniformly with derivatives in $x$ after using
endpoint-fixed local coordinates.

For the normalization, approximate $G_m$ by the cutoff function
$\mathbf1_{[\epsilon,1)}\gamma_m$ plus
$\log^{m+1}\epsilon\,\delta_0/(m+1)$, and similarly for $G_n$. On a
compact subinterval of $0<x<1$, their convolution consists of
\[
 \int_\epsilon^{x-\epsilon}\gamma_m(t)\gamma_n(x-t)\dd t
 +\int_x^1\gamma_m(t)\gamma_n(1+x-t)\dd t
\]
plus the two logarithmic counterterms multiplying $\gamma_n(x)$ and
$\gamma_m(x)$. The product of the two delta counterterms is supported
at $0$, outside that compact subinterval. Integrating the two subtracted
$a$-terms explicitly and taking the limit gives~\eqref{eq:Cmn}.
The locally uniform endpoint argument identifies this limit with the
restriction of the distributional convolution.
\end{proof}

\begin{corollary}[Four scalar identity families]\label{cor:scalar}
For $0<x<1$,
\begin{align}
 C_{0,0}(x)&=2\gamma_1(x)+\zeta(2),\label{eq:C00}\\
 C_{0,1}(x)&=\frac32\gamma_2(x)-\zeta(2)\gamma_0(x)-\zeta(3),\label{eq:C01}\\
 C_{0,2}(x)&=\frac43\gamma_3(x)-2\zeta(2)\gamma_1(x)
                    +2\zeta(3)\gamma_0(x)+2\zeta(4),\label{eq:C02}\\
 C_{1,1}(x)&=\gamma_3(x)-2\zeta(2)\gamma_1(x)
                    +2\zeta(3)\gamma_0(x)+\frac14\zeta(4).
 \label{eq:C11}
\end{align}
\end{corollary}
\begin{proof}
Restrict~\eqref{eq:product00}--\eqref{eq:product11} to $0<x<1$, where
$G_n$ agrees with $\gamma_n$ and $\Id$ agrees with the constant $-1$.
\end{proof}

\subsection{A concrete digamma integral}
Since $\gamma_0=-\psi$, the first case reads
\begin{align}
 &\int_0^x\left[\psi(t)\psi(x-t)+\psi(x)\left(\frac1t+\frac1{x-t}\right)\right]\dd t
 \label{eq:digamma-integral}\\
 &\qquad+\int_x^1\psi(t)\psi(1+x-t)\dd t-2\psi(x)\log x
 =2\gamma_1(x)+\frac{\pi^2}{6}.\notag
\end{align}
The two divergent-looking summands inside the first bracket must be
integrated together; their sum is integrable. Equation~\eqref{eq:digamma-integral}
is not an assertion that either unregularized convolution integral exists.

At $x=1/2$, the Hurwitz identity
$\zeta(s,1/2)=(2^s-1)\zeta(s)$ implies
\[
 \gamma_1(1/2)=\gamma_1-2\gamma\log2-\log^22.
\]
Thus the left side of~\eqref{eq:digamma-integral} at $1/2$ equals
\begin{equation}\label{eq:half}
 2\gamma_1-4\gamma\log2-2\log^22+\frac{\pi^2}{6}.
\end{equation}
The same method at any rational argument turns the full multiplication
table into exact relations among rational-argument Stieltjes functions,
with no integer-relation search.

\section{Differentiation and the contact-term correction}
\label{sec:contact}
For $k\ge1$, the ordinary function $\gamma_n^{(k)}(x)$ has a power-log
singularity of order $k+1$. Its canonical Hadamard finite part on $\T$
is defined by
\begin{equation}\label{eq:Fdefinition}
 \langle F_{n,k},\phi\rangle
 =\CT_{\epsilon\downarrow0}
       \int_\epsilon^1\gamma_n^{(k)}(x)\phi(x)\dd x.
\end{equation}
Here $\CT$ means the coefficient independent of $\epsilon$ after the
finite expansion in negative powers of $\epsilon$ and nonconstant powers
of $\log\epsilon$ is removed. The remaining error tends to zero.
This convention uses the coordinate $x\in(0,1)$ and no finite counterterm.
It agrees with Definition~\ref{def:G} when $k=0$.

\begin{lemma}[Higher Hurwitz poles as distributions]\label{lem:higherpoles}
For each $k\ge1$, as $z\to0$,
\begin{equation}\label{eq:higherpole}
 Z(1+k-z)=\frac{(-1)^k}{k!}\frac{\delta_0^{(k)}}z
 +\sum_{j\ge0}\frac{(-1)^jz^j}{j!}
           \Fp_{\T}[\zeta^{(j)}(k+1,x)].
\end{equation}
The finite parts in this formula have the cutoff convention
\eqref{eq:Fdefinition}.
\end{lemma}
\begin{proof}
Split $\zeta(1+k-z,x)=x^{-k-1+z}+\zeta(1+k-z,1+x)$.
Subtract the Taylor polynomial of $\phi$ at zero through degree $k$.
Its remainder gives an analytic integral near $z=0$. The subtracted
monomials give terms
$\phi^{(r)}(0)/(r!(z+r-k))$, $0\le r\le k$; only $r=k$ has a pole
at zero. Since $\langle\delta_0^{(k)},\phi\rangle=(-1)^k\phi^{(k)}(0)$,
its residue is the one stated. Expanding each regular term and each
analytic remainder is exactly the constant-term prescription for the
cutoff integrals. This proves the entire Laurent expansion.
\end{proof}

\begin{theorem}[All-index differentiation anomaly]\label{thm:anomaly}
For $n\ge0$ and $k\ge1$,
\begin{equation}\label{eq:anomaly}
 D^kG_n=F_{n,k}+(-1)^{n+1}n!e_{n+1}(k)\delta_0^{(k)}.
\end{equation}
In particular, there is no anomaly when $n\ge k$.
\end{theorem}
\begin{proof}
The meromorphic distribution family satisfies
\begin{equation}\label{eq:DZ}
 D^kZ(s)=(-1)^k(s)_kZ(s+k).
\end{equation}
This follows directly from its Fourier coefficients and the Gamma
recurrence; it also follows by repeated integration by parts in a
left half-plane and meromorphic continuation. Put $s=1-z$ and expand
\[
 (1-z)_k=k!\prod_{r=1}^k(1-z/r)
        =k!\sum_{i=0}^k(-1)^ie_i(k)z^i.
\]
On the left of~\eqref{eq:DZ}, the pole is $\delta_0^{(k)}/z$ and the
regular coefficient of $z^n$ is $D^kG_n/n!$. On the right, multiplying
the regular part of Lemma~\ref{lem:higherpoles} by the polynomial gives
$F_{n,k}/n!$: this is the finite part of the ordinary pointwise
parameter-derivative identity. Multiplying its polar part by that same
polynomial gives, in addition, the regular coefficient
$(-1)^{n+1}e_{n+1}(k)\delta_0^{(k)}$. Comparison proves the formula.
\end{proof}

The first examples are
\begin{align}
 DG_0&=F_{0,1}-\delta_0',\label{eq:anomalylow}\\
 D^2G_0&=F_{0,2}-\tfrac32\delta_0'',\notag\\
 D^2G_1&=F_{1,2}+\tfrac12\delta_0''.\notag
\end{align}
The point-supported terms disappear upon restriction to $(0,1)$, so
ordinary differentiation there cannot detect them. But convolving
$\delta_0^{(k)}$ with another distribution differentiates that other
factor everywhere. Therefore these terms cannot be discarded before
convolution, even when the desired final identity is off the singular
point.

\subsection{Relationship with the manuscript's master identity}
The pointwise formula used in the existing derivative chapter is
\begin{equation}\label{eq:pointwisemaster}
 \gamma_n^{(k)}(x)
 =(-1)^{n+k}k!n!\sum_{i+j=n}e_i(k)
              \frac{\zeta^{(j)}(k+1,x)}{j!}.
\end{equation}
It follows by expanding
$\partial_x^k\zeta(s,x)=(-1)^k(s)_k\zeta(s+k,x)$ at $s=1$,
and is consistent with Coffey's all-order formulas~\cite{Coffey}.
Theorem~\ref{thm:anomaly} neither contradicts nor modifies
\eqref{eq:pointwisemaster}; it corrects the \emph{additional, unjustified}
step of commuting canonical finite part with periodic differentiation.
This distinction should be retained in any integration into the manuscript.

\subsection{Convolutions of arbitrary differentiated Stieltjes functions}
The same argument gives a closed law for all derivative indices, not
only the zeroth Stieltjes function. Put $F_{n,0}=G_n$ and define
\[
 a_{n,k}=(-1)^{n+1}n!e_{n+1}(k)\quad(k\ge1),\qquad a_{n,0}=0.
\]
Thus $F_{n,k}=D^kG_n-a_{n,k}\delta_0^{(k)}$ when $k\ge1$.

\begin{corollary}[The full differentiated convolution calculus]
\label{cor:fulltower}
Suppose $n,m,k,l\ge0$, $r=k+l\ge1$, and the exact rule from
Theorem~\ref{thm:product} is
\[
 G_n*G_m=c_{-1}\Id+\sum_{j=0}^{n+m+1}c_jG_j.
\]
Then
\begin{align}
 F_{n,k}*F_{m,l}={}&\sum_{j=0}^{n+m+1}c_jD^rG_j
             -a_{n,k}D^rG_m-a_{m,l}D^rG_n\label{eq:fulltower}\\
 &+(c_{-1}+a_{n,k}a_{m,l})\delta_0^{(r)}.\notag
\end{align}
Away from the integer point this reduces to
\[
 \sum_{j=0}^{n+m+1}c_j\gamma_j^{(r)}(x)
          -a_{n,k}\gamma_m^{(r)}(x)-a_{m,l}\gamma_n^{(r)}(x).
\]
In particular, every such convolution has a finite reduction to
$\zeta^{(j)}(r+1,x)$ with $0\le j\le n+m+1$.
\end{corollary}
\begin{proof}
Multiply the two identities $F=D G-a\delta$ at their indicated
orders, interpreting the zero-order case through $a_{n,0}=0$.
The product of the first terms is $D^r(G_n*G_m)$, and
$D^r\Id=\delta_0^{(r)}$. Restricting off zero and applying
\eqref{eq:pointwisemaster} proves the last assertion.
\end{proof}

As an explicit higher-index example, the contact term vanishes for
$F_{1,1}=DG_1$. Hence
\begin{align}
 (F_{1,1}*F_{1,1})(x)
 ={}&-2\zeta'''(3,x)-9\zeta''(3,x)\label{eq:gamma1primepair}\\
 &+[4\zeta(2)-6]\zeta'(3,x)
   +[6\zeta(2)+4\zeta(3)]\zeta(3,x).\notag
\end{align}
This follows by twice differentiating~\eqref{eq:product11} and using
\eqref{eq:pointwisemaster}. It illustrates that the full tower still
closes at one spectral integer even when higher Stieltjes indices occur.

\section{A two-jet collapse for every polygamma convolution}
\label{sec:poly}
Let $H_0=0$ and $H_k=\sum_{j=1}^k1/j$ for $k\ge1$. Define
\[
 \Psi_k=\Fp_{\T}\psi^{(k)}\qquad(k\ge0)
\]
using the same canonical cutoff convention. Since $\gamma_0=-\psi$,
Theorem~\ref{thm:anomaly} gives
\begin{equation}\label{eq:PsiD}
 \Psi_0=-G_0,\qquad
 \Psi_k=-D^kG_0-H_k\delta_0^{(k)}\quad(k\ge1).
\end{equation}

\begin{theorem}[Polygamma convolution reduction]\label{thm:polygamma}
Let $k,l\ge0$ and $r=k+l\ge1$. Then
\begin{align}
 \Psi_k*\Psi_l={}&2D^rG_1+(H_k+H_l)D^rG_0\label{eq:polyfull}\\
 &+[H_kH_l-\zeta(2)]\delta_0^{(r)}.\notag
\end{align}
Consequently, for $0<x<1$,
\begin{equation}\label{eq:polyscalar}
 (\Psi_k*\Psi_l)(x)=(-1)^{r+1}r!
 \left[2\zeta'(r+1,x)+(2H_r-H_k-H_l)\zeta(r+1,x)\right].
\end{equation}
For $k=l=0$, the appropriate separate formula is
$\Psi_0*\Psi_0=2G_1-\zeta(2)\Id$.
\end{theorem}
\begin{proof}
Convolving~\eqref{eq:PsiD} and using that derivatives may be placed on
either factor gives
\[
 \Psi_k*\Psi_l=D^r(G_0*G_0)
 +(H_k+H_l)D^rG_0+H_kH_l\delta_0^{(r)}.
\]
The same expression is valid when one index is zero, since its harmonic
number vanishes. Apply~\eqref{eq:product00} and
$D^r\Id=\delta_0^{(r)}$. Restriction to $0<x<1$ gives
$2\gamma_1^{(r)}+(H_k+H_l)\gamma_0^{(r)}$.
Now use~\eqref{eq:pointwisemaster} with $n=0,1$:
\[
 \gamma_1^{(r)}=(-1)^{r+1}r![\zeta'(r+1,x)+H_r\zeta(r+1,x)],
 \quad \gamma_0^{(r)}=(-1)^rr!\zeta(r+1,x).
\]
This proves~\eqref{eq:polyscalar}.
\end{proof}

For example,
\begin{align}
 (\Psi_0*\Psi_1)(x)&=2\zeta'(2,x)+\zeta(2,x),\label{eq:polyexamples}\\
 (\Psi_1*\Psi_1)(x)&=-4\zeta'(3,x)-2\zeta(3,x),\notag\\
 (\Psi_0*\Psi_2)(x)&=-4\zeta'(3,x)-3\zeta(3,x).\notag
\end{align}
Although the last two convolutions have the same total derivative order,
they have different coefficients. Confusing $\Psi_k$ with the $k$th
derivative of $\Psi_0$ would incorrectly erase this distinction.

\subsection{A convergent trigamma identity with every subtraction visible}
For $0<x<1$, set $h=x/2$ and $f(t)=\psi'(t)$. Then
\begin{align}
 \mathcal T(x)={}&2\int_0^h\left[
       \frac{f(x-t)-f(x)+tf'(x)}{t^2}
        +\psi'(1+t)f(x-t)\right]\dd t\label{eq:trigamma-integral}\\
 &+\int_x^1f(t)f(1+x-t)\dd t
      -\frac{2f(x)}h-2f'(x)\log h\notag
\end{align}
satisfies
\begin{equation}\label{eq:trigamma-value}
 \mathcal T(x)=-4\zeta'(3,x)-2\zeta(3,x).
\end{equation}
Indeed, the shift relation $\psi'(t)=t^{-2}+\psi'(1+t)$ and symmetry
about $t=x/2$ reduce the first singular convolution integral to twice
its left half. Subtracting the constant and linear Taylor terms of
$f(x-t)$ leaves a bounded quotient. The constant terms of
$\int_\epsilon^ht^{-2}\dd t$ and
$\int_\epsilon^ht^{-1}\dd t$ are $-1/h$ and $\log h$,
respectively. This proves that~\eqref{eq:trigamma-integral} is exactly
the canonical finite-part convolution, so Theorem~\ref{thm:polygamma}
applies. Every integral in~\eqref{eq:trigamma-integral} is ordinary and
absolutely convergent.

\section{Antiderivatives and log-gamma convolution powers}
\label{sec:gamma}
At the integrated end no finite part is needed. Put
\begin{equation}\label{eq:gr}
 L=\log(2\pi),\qquad Z_j(x)=\zeta^{(j)}(0,x),\qquad
 g(x)=Z_1(x)=\log\Gamma(x)-\frac L2.
\end{equation}
Each $Z_j$ has mean zero. It is locally a polynomial in $\log x$ plus
a smooth function near zero, and hence belongs to every finite $L^p(0,1)$.

\begin{proposition}[Centered Stieltjes primitives]\label{prop:primitive}
For $n\ge0$, the unique mean-zero periodic distribution primitive of
$G_n$ is the integrable function
\begin{equation}\label{eq:primitive}
 A_n(x)=\frac{(-1)^{n+1}}{n+1}\zeta^{(n+1)}(0,x),\qquad DA_n=G_n.
\end{equation}
Its ordinary derivative on $(0,1)$ is $\gamma_n(x)$. The repository's
primitive based at $1$ is obtained by subtracting the value at $1$:
\[
 \int_1^x\gamma_n(t)\dd t=A_n(x)-A_n(1).
\]
\end{proposition}
\begin{proof}
Differentiate the meromorphic identity $DZ(s)=-sZ(s+1)$ at $s=0$,
using~\eqref{eq:ZLaurent}. Its coefficient of $s^{n+1}$ gives
$D\zeta^{(n+1)}(0,x)=(-1)^{n+1}(n+1)G_n$. The mean is zero because
$Z(s)$ has mean zero. Finally, a periodic distribution annihilated by
$D$ has only a zeroth Fourier coefficient, so imposing mean zero makes
the primitive unique.
\end{proof}

\begin{theorem}[All mixed integrated jets]\label{thm:integrated}
Let $d\ge2$ and $r_1,\ldots,r_d\ge0$ be integers. Write
$S=s_1+\cdots+s_d$. Then, as an identity of continuous functions after
circular convolution,
\begin{equation}\label{eq:mixedintegrated}
 Z_{r_1}*\cdots*Z_{r_d}(x)
 =\left.
 \partial_{s_1}^{r_1}\cdots\partial_{s_d}^{r_d}
 \left[\frac{\prod_{i=1}^d\Gamma(1-s_i)}{\Gamma(d-S)}
              \zeta(1-d+S,x)\right]\right|_{s_1=\cdots=s_d=0}.
\end{equation}
It holds for $0<x<1$ and at the integer point by continuous extension.
Thus any mixed convolution of the centered primitives $A_n$ reduces to
finitely many jets at the single integer $1-d$.
\end{theorem}
\begin{proof}
Before differentiation, write each Hurwitz factor as
$\zeta(s_i,x)=\Gamma(1-s_i)R_{1-s_i}(x)$ and apply
Theorem~\ref{thm:semigroup} repeatedly. For parameters sufficiently close
to zero, each factor is holomorphic as an $L^2$-valued family: the
endpoint singularities are bounded by
$x^{-\eta}(1+|\log x|^M)$ with $2\eta<1$ on each fixed derivative
order and compact parameter set. Convolution of two $L^2$ functions is
continuous, by Cauchy--Schwarz and translation continuity; further
convolution with $L^1$ factors preserves continuity. These same estimates
justify differentiating all the convolutions. The right side is
continuous at the integer point since $\Rea(1-d+S)<0$ in a small
parameter neighborhood. This proves the formula and its endpoint
extension.
\end{proof}

The case with every $r_i=1$ admits a compact all-$d$ formula. Define
$H_m^{(j)}=\sum_{a=1}^ma^{-j}$, and, for fixed $d\ge2$, set
\begin{equation}\label{eq:bd}
 b_1=H_{d-1},\qquad
 b_j=(j-1)!\big[H_{d-1}^{(j)}-\zeta(j)\big]\quad(j\ge2).
\end{equation}

\begin{theorem}[All convolution powers of centered log-gamma]\label{thm:gammapowers}
For every integer $d\ge2$,
\begin{equation}\label{eq:gammapowers}
 g^{*d}(x)=\frac1{(d-1)!}\sum_{j=0}^d\binom dj
       \cB_{d-j}(b_1,\ldots,b_{d-j})\zeta^{(j)}(1-d,x),
\end{equation}
where $\cB_0=1$ and the $b_j$ are scalar quantities in~\eqref{eq:bd}.
\end{theorem}
\begin{proof}
In~\eqref{eq:mixedintegrated}, differentiate once in each $s_i$.
Only the squarefree coefficient $s_1\cdots s_d$ matters. Modulo
$s_i^2=0$, one has
\[
 \prod_{i=1}^d\Gamma(1-s_i)
 =\prod_{i=1}^d(1+\gamma s_i)
 =\exp(\gamma S).
\]
Consequently the required mixed derivative equals the ordinary $d$th
derivative at $S=0$ of
$e^{\gamma S}\zeta(1-d+S,x)/\Gamma(d-S)$. Now
\begin{align*}
 \log\frac{e^{\gamma S}\Gamma(d)}{\Gamma(d-S)}
 &=H_{d-1}S+
   \sum_{j\ge2}\frac{H_{d-1}^{(j)}-\zeta(j)}jS^j.
\end{align*}
This follows directly from
$\Gamma(d-S)=\Gamma(1-S)\prod_{a=1}^{d-1}(a-S)$ and the Taylor series
of $\log\Gamma(1-S)$. Apply~\eqref{eq:Belldef} and the Leibniz rule;
$\Gamma(d)=(d-1)!$ supplies the prefactor.
\end{proof}

The first two nontrivial cases are
\begin{align}
 g*g={}&\zeta''(-1,x)+2\zeta'(-1,x)
                   +[2-\zeta(2)]\zeta(-1,x),\label{eq:g2}\\
 g^{*3}={}&\frac12\zeta'''(-2,x)+\frac94\zeta''(-2,x)
       +\left[\frac{21}4-\frac32\zeta(2)\right]\zeta'(-2,x)
       \label{eq:g3}\\
 &+\left[\frac{45}8-\frac94\zeta(2)-\zeta(3)\right]\zeta(-2,x).
 \notag
\end{align}
The first left side is explicitly
\[
 \int_0^xg(t)g(x-t)\dd t+\int_x^1g(t)g(1+x-t)\dd t.
\]
For uncentered $f=\log\Gamma$, mean-zero convolution annihilates all
mixed constant/nonconstant terms, so
\begin{equation}\label{eq:uncentered}
 f^{*d}(x)=\left(\frac L2\right)^d+g^{*d}(x).
\end{equation}

\begin{remark}[Convolution powers are not pointwise moments]
The value $g^{*3}(0)$ is a two-dimensional circular-convolution integral.
It is not $\int_0^1g(t)^3\dd t$. Thus~\eqref{eq:g3} does not eliminate
the Tornheim derivative coordinates of the ordinary cubic log-gamma
moment discussed in the repository. Likewise, the mixed reflected
pointwise moments in \texttt{07-reflected-moments.tex} are different
objects. This distinction prevents a spurious claim of resolving that
separate reduction problem.
\end{remark}

\section{Reflection, polylogarithmic order derivatives, and shifted correlations}
\label{sec:polylogs}
For a distribution $T$, let $\widetilde T(x)=T(-x)$; there is no complex
conjugation in this definition. Reflection changes the sign of its
Fourier index. This produces a different, polylogarithmic, multiplication
law from the unreflected Hurwitz semigroup.

\begin{theorem}[Reflected spectral master identity]\label{thm:reflected}
For $\Rea\alpha,\Rea\beta>1$ and $0\le x\le1$,
\begin{align}
 (R_\alpha*\widetilde R_\beta)(x)
 =(2\pi)^{-\alpha-\beta}\big[&e^{\pi i(\beta-\alpha)/2}
                  \Li_{\alpha+\beta}(e^{2\pi ix})\label{eq:reflected}\\
 &+e^{-\pi i(\beta-\alpha)/2}
                  \Li_{\alpha+\beta}(e^{-2\pi ix})\big].\notag
\end{align}
The right side, interpreted by its periodic Fourier coefficients,
continues to all $(\alpha,\beta)\in\C^2$ as a distribution-valued
identity. Arbitrary parameter derivatives of this equation are valid.
\end{theorem}
\begin{proof}
For $k>0$, the Fourier multiplier on the left is
\[
 (2\pi i k)^{-\alpha}(-2\pi i k)^{-\beta}
 =(2\pi k)^{-\alpha-\beta}e^{\pi i(\beta-\alpha)/2}.
\]
For $k<0$ the phase is reversed. Summing the two absolutely convergent
series proves the ordinary identity. For general parameters, the
coefficients and their parameter derivatives are polynomially bounded
on compact sets, so the paired Fourier series define the entire
continuation and justify differentiation.
\end{proof}

We write
\[
 \Li_s^{[j]}(z)=\frac{\partial^j}{\partial s^j}\Li_s(z)
\]
for an \emph{order} derivative. This must not be confused with
$z$-differentiation, which instead lowers the order through
$z\partial_z\Li_s(z)=\Li_{s-1}(z)$.

\begin{corollary}[Full shifted log-gamma correlation]\label{cor:correlation}
Put $L=\log(2\pi)$ and $A=\gamma+L$. For $0\le x\le1$, with endpoint
values understood by continuity,
\begin{align}
 &\int_0^1\log\Gamma(t)\log\Gamma(\{t-x\})\dd t
   =\frac{L^2}{4}\label{eq:correlation}\\
 &\quad+\frac1{2\pi^2}\Rea\left[
     \Li_2^{[2]}(e^{2\pi ix})-2A\Li_2^{[1]}(e^{2\pi ix})
       +\left(A^2+\frac{\pi^2}{4}\right)\Li_2(e^{2\pi ix})\right].\notag
\end{align}
Every integral here is an ordinary absolutely convergent integral.
\end{corollary}
\begin{proof}
Differentiating the Hurwitz Fourier coefficient at $s=0$ gives, for
$k\ne0$,
\[
 \widehat g(k)=\frac{\gamma+\lambda_k}{2\pi i k}.
\]
Since $g$ is real-valued, its correlation has coefficient
$\widehat g(k)\widehat g(-k)=|\widehat g(k)|^2$. For $k>0$ this is
\[
 \frac{(A+\log k)^2+\pi^2/4}{4\pi^2k^2}.
\]
The coefficients are absolutely summable. Combining the positive and
negative indices and using
$\Li_s^{[j]}(z)=\sum_{k\ge1}(-\log k)^jz^k/k^s$ near $s=2$
proves the formula for centered $g$. Restoring the mean gives $L^2/4$.
The convolution of two $L^2$ functions is continuous; alternatively,
the displayed absolutely convergent Fourier series is uniform in $x$.
At $x=0$, the integrand has only an integrable squared logarithm.
\end{proof}

At $x=0$, this recovers the classical second moment studied by
Espinosa and Moll~\cite{EM1} and already present in the manuscript:
\begin{align}
 \int_0^1\log^2\Gamma(t)\dd t
 ={}&\frac{\gamma^2}{12}+\frac{\pi^2}{48}
      +\frac{\gamma L}{6}+\frac{L^2}{3}\label{eq:secondmoment}\\
 &-\frac{\gamma+L}{\pi^2}\zeta'(2)
      +\frac{\zeta''(2)}{2\pi^2}.\notag
\end{align}
This specialization is a consistency check, not a new evaluation claim.
The full shift retains the polylogarithmic phase.

At the half shift, put $\ell=\log2$. Using
$\Li_s(-1)=(2^{1-s}-1)\zeta(s)$ in~\eqref{eq:correlation} yields
\begin{align}
 \int_0^1\log\Gamma(t)\log\Gamma(\{t-\tfrac12\})\dd t
 ={}&\frac{L^2}4-\frac{\zeta''(2)}{4\pi^2}
       +\frac{A-\ell}{2\pi^2}\zeta'(2)\label{eq:halfcorr}\\
 &+\frac{\zeta(2)}{4\pi^2}
       \left[\ell^2+2A\ell-A^2-\frac{\pi^2}4\right].\notag
\end{align}
Other rational shifts give finite Hurwitz-jet grids, as follows.

\begin{proposition}[Cyclotomic jet conversion]\label{prop:cyclotomic}
Let $q\ge1$, $p\in\Z$, and $N\ge2$ be integers. For $r\ge0$,
\begin{equation}\label{eq:cyclotomic}
 \Li_N^{[r]}(e^{2\pi ip/q})
 =q^{-N}\sum_{a=1}^q e^{2\pi ipa/q}
     \sum_{j=0}^r\binom rj(-\log q)^{r-j}\zeta^{(j)}(N,a/q).
\end{equation}
In particular, every rational-shift correlation in
Corollary~\ref{cor:correlation} reduces to finitely many jets at $s=2$.
\end{proposition}
\begin{proof}
For $\Rea s>1$, split the defining Dirichlet series for the polylogarithm
into residue classes modulo $q$:
\[
 \Li_s(e^{2\pi ip/q})
 =q^{-s}\sum_{a=1}^q e^{2\pi ipa/q}\zeta(s,a/q).
\]
Absolute locally uniform convergence permits $r$ derivatives in $s$.
The Leibniz rule gives~\eqref{eq:cyclotomic}. Its underlying distribution
and polylogarithm identities are classical~\cite{DLMFHz,DLMFLi}; here it
is the explicit conversion rule for the new convolution calculations.
\end{proof}

\section{Audit findings and proposed manuscript integration}
\label{sec:audit}
No incorrect central theorem was established in the inspected integration
and differentiation chapters. The following distinctions and localized
corrections should accompany any integration of this continuation.

\paragraph{The singular normalization must be stated.}
Equations~\eqref{eq:C00} and~\eqref{eq:polyscalar} require the canonical
finite parts in Definitions~\ref{def:G} and~\eqref{eq:Fdefinition}.
For example, the prescription $D^k\Psi_0$ is a different extension from
$\Psi_k$ when $k\ge1$. Omitting the anomaly would produce
$-4\zeta'(3,x)-6\zeta(3,x)$ instead of~\eqref{eq:intro-tri} for the
trigamma convolution. This is a correction to a tempting extension of
the existing calculus, not a claim that the manuscript prints that false
trigamma formula.

\paragraph{A local integration-by-parts clarification.}
The proof of the Hurwitz moment transform in
\texttt{07-integration.tex}, at label
\texttt{stieltjes:thm:transform}, treats the lower endpoint through the
factor $x^m$. For $m\ge1$ its vanishing is appropriate. For the base
case $m=0$, the conclusion $I_0(s)=0$ instead follows because the
primitive $\zeta(s-1,x)$ has \emph{equal} endpoint limits when
$\Rea s<1$; those limits are generally nonzero. The theorem is correct.
The proposed clarification is to state the cancellation separately rather
than allow the vanishing-endpoint explanation to cover $m=0$.

\paragraph{A PSLQ wording refinement.}
The early negative-search note at
\texttt{integral:neg:psim2} says that basis atoms must be
$\Q$-independent. Complete arithmetic independence is generally unknown
and is not an operational prerequisite for a valid exact relation.
A more precise instruction is to remove \emph{known} dependencies,
require a nonzero target coefficient, and independently prove any proposed
relation. The later discovery chapter already expresses this caution
more accurately. This is an editorial refinement, not a rejected
mathematical identity.

\paragraph{Do not reopen settled items or conflate integral types.}
At the pinned snapshot, the surviving-programme chapter marks $S_4$
proved and the cubic log-gamma moment represented through Tornheim
coordinates. The present convolution powers do not resolve its further
reduction to a smaller algebra. The finite-part closure also proves no
independence among numerical Stieltjes values at a fixed argument.

\paragraph{Placement.}
A suitable new thematic report is
\nolinkurl{Analysis/Polylogarithms/docs/reports/stieltjes-convolution/}.
The integration chapter can cite Sections~\ref{sec:gamma}--\ref{sec:polylogs};
the derivative chapter can cite Sections~\ref{sec:contact}--\ref{sec:poly}.
The full finite-part and Fourier conventions should be retained in a
standalone subsection rather than scattered among scalar tables. No
remote repository files were changed by this delivery.

\section{Verification and reproducibility}
\label{sec:verification}
The analytic proofs above are independent of numerical experiments.
The accompanying program performs two deliberately different kinds of
checks.

\subsection{Finite exact algebra}
The exact suite uses rational arithmetic and formal symbols $z_j$ for
$\zeta(j)$. It compares an exponential-series recurrence with an
independent complete Bell-polynomial recurrence; compares the product
of two one-variable normal forms with a separately assembled bivariate
Gamma quotient; checks the leading coefficient for all $0\le m,n\le4$;
checks the four printed low-index laws, applying the even-zeta relation
only where needed; compares elementary symmetric coefficients with
unsigned Stirling numbers through $k=12$; directly differentiates the
singular term $\log^n(x)/x$ through $k=6,n=6$; and checks the
squarefree-derivative coefficients of log-gamma convolution powers
through $d=7$.

There are 229 exact assertions. The emitted
\texttt{identity\_catalog.json} contains the 25 finite Stieltjes product
rules and the six log-gamma convolution-power coefficient lists. The
finite checks certify those polynomial identities relative to their
stated input expansions. They are not a proof-assistant verification of
the infinite-dimensional analytic results.

\subsection{Independent numerical diagnostics}
The numerical suite compares ordinary quadrature of the endpoint-subtracted
digamma, higher Stieltjes, and trigamma integrals with independent Stieltjes and Hurwitz-zeta
routines at $x=0.27,0.5,0.81$. It also checks the ordinary log-gamma
convolution, a complex-order semigroup instance, shifted correlations,
cyclotomic order jets, a Fourier-mode contact term, and the half-point
Stieltjes specialization. There are 31 numerical checks, including seven cross-checks of the Taylor evaluator used for $\gamma_1$. The machine-readable output records the actual precision and residual for every check. These are floating-point
diagnostics, not interval enclosures or equality certificates.

Two numerical details are important. Near a subtracted endpoint, the
code evaluates removable quotients by Taylor polynomials rather than
subtracting nearly equal floating-point numbers. For order derivatives
of polylogarithms near an integer order, a fixed small differentiation
step is used when necessary; a naively tiny default step can amplify
internal cancellation. Exact roots such as $-1$ are supplied exactly,
not as rounded complex exponentials. This is a numerical implementation
issue, not a change to the analytic formulas.

\subsection{Rebuilding}
The package uses Python with \texttt{mpmath} and \texttt{sympy}, and a
standard \LaTeX\ installation. The verification code makes no network
requests. Run
\begin{verbatim}
python verification/verify.py --exact-only
python verification/verify.py --numeric-only --dps 40
pdflatex -interaction=nonstopmode -halt-on-error article.tex
pdflatex -interaction=nonstopmode -halt-on-error article.tex
\end{verbatim}
The README gives the tested dependency versions and the contents of every
auxiliary artifact. Test output is written beside the verification
script; rerun on a working copy when preserving original diagnostics.

\section{Further research questions}
\label{sec:future}
The present work proves exact closure for a specific convolution geometry.
Several next questions have clear mathematical targets rather than merely
larger numerical searches.

\subsection{Twists, characters, and alternate endpoint normalizations}
For a nonintegral twist $\theta$, the multipliers
$(2\pi i(k+\theta))^{-\alpha}$ give a twisted spectral calculus with
no vanishing frequency $k+\theta$. Determine the analogue of the finite-part
Stieltjes normal form and identify precisely which Dirichlet or Lerch
jets replace the ordinary zeta coefficients. A new result should include its singular normalization, not merely
a nonsingular convolution identity.
A useful first target is an explicit all-index multiplication table for
the half-twist, including its point-supported counterterms.

\subsection{Affine arguments and the boundary of depth-one reduction}
The complete collapse in Theorem~\ref{thm:integrated} relies on matching
Fourier indices. If a factor has a different affine frequency, or a
pointwise product occurs before convolution, multiple independent
summation indices appear. Classify the integer-affine configurations
that still reduce to a single Hurwitz jet. For the remaining
configurations, identify the minimal natural multiple-zeta, colored
polylogarithm, or Tornheim family. This would sharpen the distinction
between convolution powers and the ordinary cubic moment rather than
blur it.

\subsection{Collision limits and renormalization compatibility}
The scalar formulas in Section~\ref{sec:integrals} are stated for
$0<x<1$. When $x\to0$, the two singularities collide. Determine a
complete local expansion that simultaneously tracks logarithmic terms,
power terms, and contact distributions, and compare taking the
finite part in $x$ with first taking finite parts in the integration
variable. A precise target is an explicit commutator formula for these
two operations at arbitrary Stieltjes indices. The two-jet polygamma
collapse gives a tractable test family.

\subsection{Cyclotomic integration certificates}
For a fixed denominator $q$, combine~\eqref{eq:cyclotomic} with the
repository's character-coordinate and distribution certificates. The
objective is a division-controlled exact certificate translating a
rational-shift correlation into the chosen character-jet basis, with
all conductor factors visible. Any rank statement should be about its
specified presentation; arithmetic independence after evaluation is a
separate problem.

\subsection{Negative-integer jets and arithmetic reductions}
Formula~\eqref{eq:gammapowers} organizes a large family of ordinary
integrals at one negative spectral integer. Which rational-argument
combinations simplify further under reflection, distribution, and
Dirichlet functional equations? The cases $d=3,4$ at denominators $3,4,5$
are concrete starting points. A meaningful reduction claim must specify
the comparison algebra; replacing a Hurwitz jet by an equivalent
Dirichlet derivative is a coordinate change unless that target algebra
is explicitly part of the objective.

\subsection{Minimal formalization interfaces}
A proof-assistant development can begin with finite Bell recurrences,
the Gamma-quotient coefficient identity, and the elementary-symmetric
contact coefficients. Its analytic interface then needs Fourier
uniqueness for periodic distributions, convolution multiplication,
and the endpoint subtraction lemma. Establishing the distributional
Laurent expansion once should make the remaining spectral-jet proofs
short. None of those analytic formalization steps is claimed complete
in this delivery.

\section{Conclusion}
A single Fourier normalization links four apparently different layers:
ordinary log-gamma convolution integrals, generalized Stieltjes finite
parts, polygamma convolutions, and polylogarithmic correlations. The
nonsingular Hurwitz semigroup is the classical organizing identity.
At its singular parameter, the residue $\delta_0-1$ creates a closed
one-generator Stieltjes convolution algebra. Its higher poles determine
the exact derivative contact terms. Retaining those terms yields the
correct two-jet polygamma reduction; moving in the opposite direction
produces all log-gamma convolution powers and their mixed primitive
analogues. These are explicit identities with proved domains and
normalizations, rather than conjectures supported only by small
numerical residuals.

\clearpage
\begin{thebibliography}{99}
\bibitem{ProveIt}
V. Reshetnikov and ProveIt contributors,
\emph{ProveIt, Analysis/Polylogarithms: consolidated research manuscript},
snapshot \pin, 10 October 2026. Source directory:
\url{https://github.com/VladimirReshetnikov/ProveIt/tree/0efdbaa6944357e578a8839ec88d040c6e05681c/Analysis/Polylogarithms/docs/manuscript}.

\bibitem{Sun}
Zhi-Hong Sun, \emph{On the properties of invariant functions},
arXiv:2209.14625 (2022), especially Remark 3.2.
\url{https://arxiv.org/abs/2209.14625}.

\bibitem{DLMFHz}
NIST Digital Library of Mathematical Functions,
Section 25.11, \emph{Hurwitz Zeta Function}; in particular the shift,
Fourier, and parameter-derivative formulas. Accessed 10 October 2026.
\url{https://dlmf.nist.gov/25.11}.

\bibitem{DLMFLi}
NIST Digital Library of Mathematical Functions,
Section 25.12, \emph{Polylogarithms}. Accessed 10 October 2026.
\url{https://dlmf.nist.gov/25.12}.

\bibitem{EM1}
Olivier R. Espinosa and Victor H. Moll,
\emph{On some definite integrals involving the Hurwitz zeta function},
arXiv:math/0012078 (2000).
\url{https://arxiv.org/abs/math/0012078}.

\bibitem{EM2}
Olivier R. Espinosa and Victor H. Moll,
\emph{On some integrals involving the Hurwitz zeta function: part 2},
arXiv:math/0107082 (2001).
\url{https://arxiv.org/abs/math/0107082}.

\bibitem{Coffey}
Mark W. Coffey, \emph{Series representations for the Stieltjes constants},
arXiv:0905.1111 (2009), including the all-order parameter-derivative
formulas. \url{https://arxiv.org/abs/0905.1111}.
\end{thebibliography}
\end{document}
